% Part: lambda-calculus
% Chapter: syntax
% Section: eta

\documentclass[../../../include/open-logic-section]{subfiles}

\begin{document}

\olfileid{lam}{syn}{eta}
\olsection{Konversi-$\eta$}

Ana relasi liya ing suku-$\lambda$.
Ing \olref[fv]{sec} kita nggunakake
conto $\lambd[x][(fx)]$, kang
nampa argumen banjur ngetrapake
$f$ marang argumen iku.
Kanthi tembung liya, iki fungsi
kang padha karo~$f$:
$\lambd[x][(fx)]N$ lan $fN$
loro-lorone direduksi dadi~$fN$.
Kita nggunakake reduksi-$\eta$
(lan ekstensi-$\eta$) kanggo
ngandharake gagasan iki.

\begin{defn}[Kontraksi-$\eta$, $\eredone$]
  \ollabel{defn:beredone}
  \emph{Kontraksi-$\eta$} ($\eredone$)
  yaiku relasi kompatibel paling cilik
  ing suku kang nyukupi syarat iki:
  \[
    \lambd[x][M x] \eredone M \text{ yen } x \notin FV(M)
  \]
\end{defn}

\begin{defn}[Reduksi-$\beta\eta$, $\bered$] \ollabel{defn:bered}
  \emph{Reduksi-$\beta\eta$} ($\bered$)
  yaiku relasi refleksif lan transitif
  paling cilik ing suku kang ngemot
  $\bredone$ lan $\eredone$.
  Tegese, kajaba aturan refleksivitas
  lan transitivitas, ana rong aturan iki:
  \begin{enumerate}
  \item Yen $M \bredone N$, mula $M \bered N$. \ollabel{defn:bered3}
  \item Yen $M \eredone N$, mula $M \bered N$. \ollabel{defn:bered4}
  \end{enumerate}

\end{defn}

\begin{defn}
  Kita nggedhekake relasi ekuivalensi
  $\equal$ kanthi aturan konversi-$\eta$:
  \[
  \lambd[x][f x] \equal f \text{ yen } x \notin FV(f)
  \]
  lan nglambangake relasi kang
  digedhekake iku minangka
  $\equal[\eta]$.
\end{defn}

Ekuivalensi-$\eta$ wigati amarga
ana gandhengane karo ekstensionalitas
suku lambda:

\begin{defn}[Ekstensionalitas]
  Kita nggedhekake relasi ekuivalensi
  $\equal$ kanthi aturan (\ext):
  \begin{center}
  Yen $Mx \equal Nx$, mula $M \equal N$,
  kanthi syarat $x \notin FV(MN)$.
  \end{center}
  Relasi kang digedhekake iku
  kita lambangake $\equal[\ext]$.
\end{defn}

Kurang luwih, aturan iki ngandharake
yen rong suku minangka fungsi kudu
dianggep padha yen tumindake padha
marang argumen kang padha.

Saiki kita buktekake yen aturan
$\eta$ menehi ekstensionalitas
pas, tanpa nambah relasi liyane.

\begin{thm}
  $M \equal[\ext] N$ yen lan mung yen $M \equal[\eta] N$.
\end{thm}

\begin{proof}
  Sepisan, kita buktekake yen
  $\equal[\eta]$ katutup marang
  aturan ekstensionalitas. Tegese,
  aturan $\ext$ ora nambah apa-apa
  marang $\equal[\eta]$.
  Mula $\equal[\eta]$ ngemot
  $\equal[\ext]$, lan yen
  $M \equal[\ext] N$, mula
  $M \equal[\eta] N$.

  Kanggo mbuktekake katutupan
  $\equal[\eta]$ marang \ext,
  gunakake induksi ing derivasi.
  Yen langkah pungkasan aturan
  kasebut ngasilake $M
  \equal[\ext] N$ liwat \ext{},
  hipotesis
  induksi marang premisé menehi
  $Mx \equal[\eta] Nx$.
  Miturut aturan kompatibilitas
  kanggo abstraksi,
  kita oleh $\lambd[x][Mx]
  \equal[\eta] \lambd[x][Nx]$
  ing $\equal$.
  Amarga syarat kasegeran iku
  nyakup loro-lorone
  suku, konversi-$\eta$ ing
  sisih loro menehi
  $M \equal[\eta] N$.

  Semono uga, kita buktekake
  yen aturan $\eta$ wis kakandhut
  ing $\equal[\ext]$.
  Kanggo $\lambd[x][Mx]$ lan
  $M$ kanthi $x \notin FV(M)$,
  kita nduweni
  $(\lambd[x][Mx])x \equal[\ext] Mx$,
  mula kita oleh
  $\lambd[x][Mx] \equal[\ext] M$
  miturut aturan \ext{}.
\end{proof}

\end{document}
